\subsection{The original finite phase is well defined}
\label{global-gram-finite-phases}

Retain the original arithmetic setting, its level subgroup character,
its modulus $g>0$, and all its original finite vectors.  In this
subsection $\mathfrak d$ denotes the extension to $\mathcal O_E$ of
the different of $F/\mathbb Q$, the source's \texttt{ddS}; it is
not replaced by a unit ideal.  The original lattice splitting is
\begin{equation}\label{global-gram-phase-split}
 L=\bigoplus_{a=0}^{r-1}\mathfrak a_a e_a\ \oplus L_b\ \oplus
                      \bigoplus_{a=0}^{r-1}\mathfrak a_a^{-1}f_a,
 \qquad \mathfrak a_a=\mathfrak d^{a-r}.
\end{equation}
The hyperbolic vectors satisfy $h(e_a,f_b)=\delta_{ab}$, with
$h$ conjugate linear in its first argument.  The lattice $L_b$ is
orthogonal to them and contains the original $z=z_u$.
The fixed finite lattice $L_f$ is the completion of $L$ in the
original coordinates.  Write $K_0$ for its full stabilizer.

The original finite approximation and $x=x_u\in K_0$ give, at
every prime dividing $g$,
\begin{equation}\label{global-gram-phase-congruence}
 e_a\equiv e_a^0,\qquad f_a\equiv f_a^0,\qquad
 z\equiv xz^0\pmod{gL_f},\qquad xe_a=e_a,\quad xf_a=f_a.
\end{equation}
These are localized assertions.  They follow by inverting the
prescribed common integer $C$ coprime to $g$; no assertion that
the rational hyperbolic vectors are globally integral is used.
The vectors $e^0,f^0,z^0,b_i,b_i^\vee$ are in $L_f$ and satisfy
the original integral Gram congruences.  In particular, modulo
$g$, $e^0,f^0$ are dual hyperbolic frames,
$b_i,b_i^\vee$ are orthogonal to them and $z^0$, and
$h(b_i^\vee,b_j)=\delta_{ij}\epsilon$, with the original finite
idempotent $\epsilon$.\cite{global-gram-phase-source}

Set $w_a=e_{a+1}$ for $a<r-1$ and $w_{r-1}=z$, with the same
convention for $w_a^0$.  When $r=0$ all root sums are empty and
the last-index assertions below are absent.  The original two
character arguments are
\begin{align*}
 B(n)&=\sum_a h(f_a,nw_a),\\
 A(\kappa)&=\sum_a h(f_a^0,\kappa w_a^0)
                  +\sum_{i+1<k}h(b_i^\vee,\kappa b_{i+1}).
\end{align*}
Thus $\psi_{u,f}(n)=\Psi_f(B(n)/g)$ and
$\chi_2(\kappa)=\Psi_f(A(\kappa)/g)$, with the very same finite
additive character $\Psi_f$ as in the original construction.

\begin{theorem}[Exact unipotent phase matching]
\label{global-gram-phase-match}
For every $n\in N_u(\mathbb A_{F,f})\cap K_0$,
\[
 \kappa=x_u^{-1}nx_u\in K_2,
 \qquad \chi_2(\kappa)=\psi_{u,f}(n).
\]
If $n\in N_u(\mathbb A_{F,f})$, $h\in H_u(\mathbb A_{F,f})$
and $x_u^{-1}nhx_u\in K_2$, then
\[
 x_u^{-1}hx_u\in K_2,\qquad
 \chi_2(x_u^{-1}nhx_u)
       =\psi_{u,f}(n)\chi_2(x_u^{-1}hx_u).
\]
Consequently the original scalar seed formula
\eqref{global-gram-original-seed} is independent of its finite
decomposition: the original \texttt{PhWellDef} premise is supplied.
\end{theorem}

\begin{proof}
\emph{The inverse different is an actual annihilator.}
For $y\in\mathfrak d^{-1}$ and $o\in\mathcal O_E$,
$\operatorname{Tr}_{E/F}(oy)$ belongs to the inverse different
of $F/\mathbb Q$.  Indeed multiplication by every element of
that different makes $oy$ integral, so makes its relative trace
integral.  Trace duality then gives
$\operatorname{Tr}_{E/\mathbb Q}(oy)\in\mathbb Z$.
For a finite adele $v\in\widehat{\mathfrak d}^{-1}$, the finite
Chinese remainder theorem chooses $q\in\mathfrak d^{-1}$ such
that $v-q\in\widehat{\mathcal O}_E$.  The character is trivial
on the latter lattice.  Its original rational product formula
therefore gives
\[
 \Psi_f(v)=\Psi_f(q)
          =\exp(-2\pi i\operatorname{Tr}_{E/\mathbb Q}q)=1.
\]
This uses the actual extended different, including all its
ramified prime factors.

\emph{The unipotent is in the original level subgroup.}
Work first in the product of the local integer rings at primes
dividing $g$; this product is preserved by conjugation.  By
\eqref{global-gram-phase-congruence}, the actual $e_a,f_a,z$
are integral there.  For a vector in the span of the $e_a$,
its coefficient of $e_a$ is $h(f_a,\cdot)$; hence an integral
vector in that span has integral coefficients.  The same is
true for every initial subflag.  Since $n$ preserves $L_f$ and
is unipotent for the original flag, reduction modulo $g$ gives
\[
 (\kappa-1)e_a^0\in\langle e_b^0:b<a\rangle,
 \qquad (\kappa-1)z^0\in\langle e_b^0:b<r\rangle.
\]
Here and in the next display the spans are over the localized
integer ring modulo $g$.  For the third level condition put
$\beta=xb_i$.  Its pairings $h(e_a,\beta)=h(e_a,b_i)$ lie
in $g\mathcal O_{E,v}$.  The corrected vector
\[
          \widetilde\beta=\beta-\sum_a h(e_a,\beta)f_a
\]
is integral, lies in the orthogonal complement of the isotropic
span, and is congruent to $\beta$ modulo $g$.  Thus
$(n-1)\widetilde\beta$ is an integral vector in that isotropic
span.  Applying $x^{-1}$ and reducing gives
\[
               (\kappa-1)b_i\in\langle e_a^0:a<r\rangle.
\]
This is stronger than the original third level condition, which
also allows the earlier $b_j$.  At primes not dividing $g$,
$gL_f=L_f$, so all three congruence conditions are automatic.
Since $x,n\in K_0$, also $\kappa\in K_0$.  These are exactly
the original defining conditions for $K_2$.

\emph{The complete character argument.}
First $A(\kappa)\in\widehat{\mathcal O}_E$: all vectors in
its expression and their images under $\kappa$ lie in the
original integral Hermitian lattice.  Next
\begin{equation}\label{global-gram-phase-B-different}
                         B(n)\in\widehat{\mathfrak d}^{-1}.
\end{equation}
For $a<r-1$, the $e_a$ coefficient of $(n-1)e_{a+1}$ belongs
to $\widehat{\mathfrak a_a\mathfrak a_{a+1}^{-1}}
=\widehat{\mathfrak d}^{-1}$ by preservation of
\eqref{global-gram-phase-split}.  For the last term,
$(n-1)z\in L_f$ and its $e_{r-1}$ coefficient belongs to
$\widehat{\mathfrak a}_{r-1}=\widehat{\mathfrak d}^{-1}$.
The identity terms have zero pairings with the corresponding
$f_a$.  This proves \eqref{global-gram-phase-B-different}; the
argument uses the original fractional ideals, without declaring
the flag lattice free or saturated.

At primes dividing $g$, unitarity of $x$ and
\eqref{global-gram-phase-congruence} give
\[
 h(f_a^0,\kappa w_a^0)=h(xf_a^0,nxw_a^0)
                         \equiv h(f_a,nw_a)\pmod{g}.
\]
All paired vectors here are integral at these primes, and $n$
preserves the local lattice.  The third level congruence just
proved and the original Gram congruences also give
\[
 h(b_i^\vee,\kappa b_{i+1})
 \equiv h(b_i^\vee,b_{i+1})=0\pmod{g}.
\]
Therefore $(A(\kappa)-B(n))/g$ is integral at primes dividing
$g$.  At every other prime $g$ is a unit, and integrality of
$A$ together with \eqref{global-gram-phase-B-different} puts
the same difference in the inverse extended different.  Globally,
\[
          (A(\kappa)-B(n))/g\in\widehat{\mathfrak d}^{-1}.
\]
Applying its actual annihilator proves the exact phase equality.
In particular this argument keeps the contribution outside the
level primes; a congruence only modulo $g$ would not suffice.

\emph{Splitting a product in the integral stabilizer.}
Let $p$ be projection off the hyperbolic pairs onto their
orthogonal complement.  The original splitting
\eqref{global-gram-phase-split} proves $pL_f\subset L_f$.
The element $h$ fixes every $e_a,f_a$ and preserves this
orthogonal complement.  If $nh\in K_0$ and $y\in L_f$, set
$q=hpy$.  Then $nhpy\in L_f$, while $nq-q$ lies in the
isotropic span.  Hence $p(nq)=q\in L_f$.  Also
$hy=(y-py)+q\in L_f$.  Apply the same argument to
$(nh)^{-1}=(h^{-1}n^{-1}h)h^{-1}$ to prove $hL_f=L_f$.
It follows that both $h$ and $n$ lie in $K_0$.

Now $x^{-1}nx\in K_2$ by the preceding argument.  If
$x^{-1}nhx\in K_2$, the subgroup law gives
$x^{-1}hx\in K_2$, and its original character law gives the
second assertion of the theorem.

\emph{Independence of the seed decomposition.}
The functions $n\mapsto B(n)$ are additive on $N_u$ and
unchanged by $H_u$-conjugation: they are the adjacent
coefficients of its unipotent flag, and $H_u$ fixes the actual
$e_a,f_a,z$.  Thus $\psi_{u,f}$ is a character of $N_u$
invariant under $H_u$.  If
$nhxk=n'h'xk'$, put
\[
 q=h'^{-1}n'^{-1}nh',\qquad \ell=h'^{-1}h.
\]
Then $q\in N_u$ and $x^{-1}q\ell x=k'k^{-1}\in K_2$.
The just-proved product statement yields
$x^{-1}\ell x\in K_2$ and
\[
 \chi_2(k'k^{-1})
       =\psi_{u,f}(n')^{-1}\psi_{u,f}(n)\chi_2(x^{-1}\ell x).
\]
The original column law is
\[
       T(\xi)(1,h)=\chi_2(x^{-1}\ell x)T(\xi)(1,h').
\]
Substitution gives, for every $\xi$,
\[
 \psi_{u,f}(n)\chi_2(k)T(\xi)(1,h)
 =\psi_{u,f}(n')\chi_2(k')T(\xi)(1,h').
\]
Taking $\xi=W_u(g)$ proves the literal \texttt{PhWellDef}
condition.
\end{proof}

In an existing original arithmetic setting this discharges the
explicit local decomposition premise of
Theorem~\ref{global-gram-ifinite} and
Corollary~\ref{global-gram-poincare}.  It changes no lattice,
finite phase, support, or scalar normalization.  The separate
class-frame regrouping and raw-kernel identification still have
to be supplied before the scalar Gram identity can be used as
the final restricted operator identity.
